\section{Coverage inside sets of polynomial Gibbs mass}
\label{poly-sec:coverage}

The fixed-depth crossing argument in Appendix~\ref{fixed-sec:crossings}
uses polynomially many reference samples. The neighborhood bound
$e^{-D_3n^\kap}$ in Proposition~\ref{prop:coverage} is smaller than every
fixed inverse power of $n$, so it does not supply those references.
We therefore apply the same energy and rotation mechanism with a
perturbation variance of order $\log n/n$. This yields polynomial
overlap mass inside sets of polynomial Gibbs mass, uniformly over every
subset of the cube. We give the argument from the energy and rotation
estimates in Section~\ref{sec:energy}, so that the fixed-depth route can
be read independently of the growing-depth construction.

Throughout this appendix, fix $\beta>1$ and use the off-diagonal
Hamiltonian of Equation~\eqref{eq:model}. Write
$H^J=H_n^J$, $Z^J=Z_{n,\beta}^J$, and $\pi^J=\pi_{n,\beta}^J$.
For a nonempty $S\subset\Sig$, define its conditional Gibbs measure by
\[
 \pi_S^J(x)=\frac{\pi^J(x)\mathbf1_{\{x\in S\}}}{\pi^J(S)}.
\]
Every nonempty set has positive Gibbs mass. All Gaussian disorder vectors
below have independent standard normal coordinates indexed by $i<j$.

\begin{proposition}[Polynomial overlap coverage]
\label{poly-prop:coverage}
There exists $q_*=q_*(\beta)\in(0,1)$ such that, for every fixed $D_1>0$,
there exists a fixed $D_2>0$ for which
\begin{equation}\label{poly-eq:coverage-uniform}
 \PP_J\!\left(
  \text{for every }S\subset\Sig\text{ with }\pi^J(S)\ge n^{-D_1},\quad
  (\pi_S^J)^{\otimes2}\{|R|\ge q_*\}\ge n^{-D_2}
 \right)\longrightarrow1.
\end{equation}
The threshold $q_*$ does not depend on $D_1$.
\end{proposition}

The proof compares two estimates for the change in $\log Z$ under a small
Gaussian rotation. If a set has too few overlapping pairs, a fresh Gaussian
perturbation increases its restricted partition sum. The energy upper tail
limits the loss caused by decreasing the coefficient of the original
disorder. Together these give an increase that the rotation estimate makes
unlikely. We first prove the required gain for an arbitrary deterministic
probability measure, so that it can later be applied after conditioning on
$J$.

\subsection{Gain under independent Gaussian disorder}

\begin{lemma}[Gaussian gain from few overlapping pairs]
\label{poly-lem:coverage-gain}
Fix $M>0$, $q\in(0,1)$, and $D\ge\beta^2M/2+2$.
For all sufficiently large $n$, set $t=M\log n/n$.
If $\nu$ is any deterministic probability measure on $\Sig$ satisfying
\[
 \nu^{\otimes2}\{|R|\ge q\}\le4n^{-D},
\]
and $G$ is standard Gaussian, define
$W=\nu(e^{\beta\sqrt t H^G})$.
There is a constant $C_\beta$, independent of $M,q,D,n,\nu$, such that
\begin{equation}\label{poly-eq:coverage-gain}
 \PP_G\!\left(
 \log W\ge\frac{\beta^2t(n-1)}4
       -\beta^2ntq-C_\beta\sqrt{nt}-\log2
 \right)\ge\frac12.
\end{equation}
\end{lemma}

\begin{proof}
For $x\in\Sig$ let $h_x=(n^{-1/2}x_ix_j)_{i<j}$, so that
$H^G(x)=h_x\cdot G$. In particular,
\[
 |h_x|^2=\frac{n-1}{2},\qquad
 \EE_G[H^G(x)H^G(y)]=\frac{nR(x,y)^2-1}{2}.
\]
Set
\[
 m_0=\frac{\beta^2t(n-1)}4,\qquad
 A=\frac{\beta^2ntq^2}{2}.
\]
The Gaussian exponential moment formula gives $\EE_GW=e^{m_0}$.
For the second moment, split pairs according to whether $|R|<q$:
\begin{align}
 \frac{\EE_GW^2}{(\EE_GW)^2}
 &=\nu^{\otimes2}\!\left(e^{\beta^2t(nR^2-1)/2}\right)\notag\\
 &\le e^A+4n^{-D}e^{\beta^2nt/2}
 \le e^A+4n^{-2}
 \le2e^A.
 \label{poly-eq:coverage-moments}
\end{align}
The exponent assumption on $D$ gives the second inequality, and the last
one holds for $n\ge2$. Paley--Zygmund therefore implies
\begin{equation}\label{poly-eq:coverage-pz}
 \PP_G(\log W\ge m_0-\log2)\ge\frac18e^{-A}.
\end{equation}
This probability may be small. Gaussian concentration will turn it into
a lower bound of probability at least one half, while controlling the loss
in the value of $\log W$.

Write $Y(G)=\log W$. Differentiation expresses its gradient as
$\beta\sqrt t$ times a convex combination of the vectors $h_x$.
Consequently $Y$ is Lipschitz with constant at most
\[
 L_W=\beta\sqrt{\frac{t(n-1)}2}.
\]
The Gaussian moment bound in Equation~\eqref{static-eq:gaussian-mgf}
of Lemma~\ref{static-lem:gaussian-tools}, followed by Chernoff's inequality,
gives the following one-sided tails for this deterministic function of $G$:
\[
 \PP_G(Y-\EE_GY\ge s)\le e^{-s^2/(2L_W^2)},\qquad
 \PP_G(Y-\EE_GY\le-s)\le e^{-s^2/(2L_W^2)}
 \quad(s>0).
\]
If $m_0-\log2>\EE_GY$, comparison of the upper-tail bound with
Equation~\eqref{poly-eq:coverage-pz} gives
\[
 \EE_GY\ge m_0-\log2-L_W\sqrt{2(A+\log8)}.
\]
If $m_0-\log2\le\EE_GY$, the same inequality holds immediately.
Using $\sqrt{a+b}\le\sqrt a+\sqrt b$, we thus obtain
\begin{equation}\label{poly-eq:coverage-mean-gain}
 \EE_GY\ge m_0-\log2
        -\frac{\beta^2ntq}{\sqrt2}
        -\beta\sqrt{nt\log8}.
\end{equation}
The lower-tail bound also gives
$Y\ge\EE_GY-L_W\sqrt{2\log2}$ with probability at least $1/2$.
Combining this with Equation~\eqref{poly-eq:coverage-mean-gain}, and
bounding $\beta^2ntq/\sqrt2$ by $\beta^2ntq$, proves the stated estimate.
One admissible choice is
\[
 C_\beta=\beta\bigl(\sqrt{\log8}+\sqrt{\log2}\bigr).
\]
This constant is independent of the measure $\nu$ and of the parameters
$M,q,D$.
\end{proof}

\subsection{Uniform coverage over all subsets}

\begin{proof}[Proof of Proposition~\ref{poly-prop:coverage}]
Let $e_1\in(0,\beta/2)$ and $c_E>0$ be the constants from
Lemma~\ref{static-lem:coverage-energy}. Thus, outside a disorder event
of probability at most $e^{-c_En}$,
\begin{equation}\label{poly-eq:energy-event}
 \pi^J\{H^J>ne_1\}\le e^{-c_En}.
\end{equation}
Define
\[
 g=\frac{\beta^2}{4}-\frac{\beta e_1}{2}>0.
\]
Choose the overlap threshold before choosing any mass exponent:
\begin{equation}\label{poly-eq:coverage-constants}
 0<q_*<\min\left\{1,\frac{g}{8\beta^2}\right\}.
\end{equation}
It depends only on $\beta$. We choose this same threshold in
Proposition~\ref{prop:coverage}: both coverage arguments use the same
$e_1$ and permit the strict bound in
Equation~\eqref{poly-eq:coverage-constants}. Now fix $D_1>0$ and choose
\begin{equation}\label{poly-eq:coverage-parameters}
 M\ge\frac{8(D_1+1)}g,\qquad
 D_2=\frac{\beta^2M}{2}+2,\qquad
 t=\frac{M\log n}{n},\qquad \theta=\arcsin\sqrt t.
\end{equation}
For all sufficiently large $n$, $0<t<1/2$.
All constants have been fixed independently of $n$ and of the disorder.

Let $\mathcal B_n$ be the event that
Equation~\eqref{poly-eq:energy-event} holds and that some $S\subset\Sig$
satisfies
\begin{equation}\label{poly-eq:violating-set}
 \pi^J(S)\ge n^{-D_1},\qquad
 (\pi_S^J)^{\otimes2}\{|R|\ge q_*\}<n^{-D_2}.
\end{equation}
We will bound the probability of this existence event by selecting one
violating set. Fix a deterministic ordering of the subsets of the finite
cube and take its first member satisfying
Equation~\eqref{poly-eq:violating-set}. For each fixed nonempty set, both
Gibbs quantities in that equation are continuous functions of $J$.
The resulting selector $S=S(J)$ is therefore measurable on
$\mathcal B_n$. In particular, it is fixed when we condition on $J$.

For such a fixed $J$, remove the part of $S$ above the energy threshold:
\[
 T=S\cap\{x:H^J(x)\le ne_1\},\qquad \nu=\pi_T^J.
\]
For all sufficiently large $n$, uniformly over $J\in\mathcal B_n$,
\begin{equation}\label{poly-eq:trimmed-mass}
 \pi^J(T)\ge\pi^J(S)-e^{-c_En}
       \ge\frac12\pi^J(S)\ge\frac12n^{-D_1}.
\end{equation}
In particular, $T$ is nonempty. Since $T\subset S$, the two conditioning
denominators give
\begin{equation}\label{poly-eq:trimmed-pairs}
 \nu^{\otimes2}\{|R|\ge q_*\}
 \le\left(\frac{\pi^J(S)}{\pi^J(T)}\right)^2
       (\pi_S^J)^{\otimes2}\{|R|\ge q_*\}
 <4n^{-D_2}.
\end{equation}
The trimmed measure is measurable in $J$ as well. Draw $G$ independently
of $J$. Conditional on $J$, the measure $\nu$ is deterministic, so
Lemma~\ref{poly-lem:coverage-gain} applies with $q=q_*$ and $D=D_2$.
Write
\[
 W=\nu(e^{\beta\sqrt t H^G}),\qquad
 J_\theta=\sqrt{1-t}\,J+\sqrt t\,G.
\]
No part of the choice of $S,T$, or $\nu$ uses this fresh disorder.

We now compare this conditional gain with the full partition sum.
Linearity in the disorder, restriction to $T$, and
$\sqrt{1-t}-1<0$ imply
\begin{align*}
 \log\frac{Z^{J_\theta}}{Z^J}
 &=\log\pi^J\!\left(
      e^{\beta(\sqrt{1-t}-1)H^J+\beta\sqrt t H^G}\right)\\
 &\ge-D_1\log n-\log2
       -\beta(1-\sqrt{1-t})ne_1+\log W.
\end{align*}
Thus Lemma~\ref{poly-lem:coverage-gain} gives, with conditional
probability at least $1/2$,
\begin{align}
 \log\frac{Z^{J_\theta}}{Z^J}
 &\ge-D_1\log n-2\log2
       -\beta(1-\sqrt{1-t})ne_1
       +\frac{\beta^2t(n-1)}4
       -\beta^2ntq_*-C_\beta\sqrt{nt}\notag\\
 &\ge nt\left(g-\beta^2q_*-\frac{D_1}{M}
                         -\frac{\beta e_1t}{2}\right)
       -C_\beta\sqrt{nt}-2\log2-\frac{\beta^2t}{4}.
 \label{poly-eq:partition-gain}
\end{align}
The last line uses
$1-\sqrt{1-t}\le t/2+t^2/2$ for $0\le t<1$.
Our choices ensure $\beta^2q_*<g/8$ and $D_1/M<g/8$.
Since $t\to0$ and $nt=M\log n\to\infty$, the remaining errors in
Equation~\eqref{poly-eq:partition-gain} are $o(nt)$. Consequently,
for all sufficiently large $n$ and every $J\in\mathcal B_n$,
\begin{equation}\label{poly-eq:bad-increment}
 \PP_G\!\left(
   \log Z^{J_\theta}-\log Z^J\ge\frac g2nt\,\middle|\,J
 \right)\ge\frac12.
\end{equation}
The lower bound on $n$ depends only on the fixed parameters and is
uniform in $J$ and its selected set.

It remains to bound this increment under the joint law of $(J,G)$.
For $t\le1/2$, $\theta\le\sqrt{2t}$. The unconditional rotation
estimate in Lemma~\ref{static-lem:coverage-rotation} therefore yields
\begin{align*}
 \PP_{J,G}\!\left(
   \log Z^{J_\theta}-\log Z^J\ge\frac g2nt\right)
 &\le\exp\!\left(-\frac{g^2nt}{8\beta^2}\right)\\
 &=n^{-g^2M/(8\beta^2)}.
\end{align*}
Integrating Equation~\eqref{poly-eq:bad-increment} over
$\mathcal B_n$ bounds $\PP_J(\mathcal B_n)$ by twice this probability.
Adding the exceptional event from
Lemma~\ref{static-lem:coverage-energy}, we conclude that the failure
probability in Equation~\eqref{poly-eq:coverage-uniform} is at most
\begin{equation}\label{poly-eq:failure-bound}
 e^{-c_En}+2n^{-g^2M/(8\beta^2)}\longrightarrow0.
\end{equation}
The selector controls the existence of any violating subset, without a
union bound over subsets of the cube. The choice of $q_*$ in
Equation~\eqref{poly-eq:coverage-constants} preceded the choice of
$D_1$, as required.
\end{proof}

\subsection{Target coverage within arbitrary subsets}

For the crossing argument, we need a bound on the total Gibbs mass of
targets lacking polynomial Gibbs mass of overlapping references in the
same subset. Both masses below use the original Gibbs measure. This is
the form needed when targets have Gibbs marginals and reference banks
contain unconditioned Gibbs samples. Uniformity of the pair statement
gives the target bound by applying it to the set of points that would
violate it.

\begin{corollary}[Target coverage within arbitrary subsets]
\label{poly-cor:coverage-cells}
Let $q_*$ be as in Proposition~\ref{poly-prop:coverage}.
For every fixed $D_1>0$, there exists a fixed $D_3>0$ such that, with
probability tending to one over $J$, every $S\subset\Sig$ satisfies
\begin{equation}\label{poly-eq:coverage-cells}
 \pi^J\!\left(
  \left\{x\in S:
   \pi^J\{b\in S:|R(x,b)|\ge q_*\}<n^{-D_3}
  \right\}\right)<n^{-D_1}.
\end{equation}
\end{corollary}

\begin{proof}
Fix $D_1>0$, take the corresponding $D_2$ from
Proposition~\ref{poly-prop:coverage}, and choose any fixed
$D_3>D_1+D_2$. Work on the simultaneous event in that proposition.
Given $S\subset\Sig$, denote the set inside the outer measure in
Equation~\eqref{poly-eq:coverage-cells} by $B$.
If $m=\pi^J(B)\ge n^{-D_1}$, then $B$ is nonempty, and $B\subset S$
implies
\begin{align*}
 (\pi_B^J)^{\otimes2}\{|R|\ge q_*\}
 &=\frac1{m^2}\sum_{x\in B}\pi^J(x)
             \pi^J\{b\in B:|R(x,b)|\ge q_*\}\\
 &<\frac{n^{-D_3}}m
 \le n^{-D_3+D_1}<n^{-D_2}.
\end{align*}
This contradicts Proposition~\ref{poly-prop:coverage} applied to $B$.
Hence $\pi^J(B)<n^{-D_1}$. This implication is deterministic on the
same disorder event for every $S$, including cells selected after
observing $J$. For $S=\varnothing$ the assertion is immediate.
\end{proof}
